\appendix
\section{Search Strategy Validation}
\label{ap:search-validation}

This appendix documents the validation process applied to the search strategy before
executing the final queries against the databases.
The validation combines two procedures: auditing the strings against comparable systematic
reviews and checking recall by means of a seed-paper set operating as a Quasi-Gold Standard
(QGS).
The last subsection presents the full PRISMA-S report~\citep{rethlefsen2021prismas}.

\subsection{Pre-execution audit}
\label{ap:auditoria}

The pre-execution audit consulted six secondary studies published
between 2024 and 2026 on \textit{Large Language Model} (LLM)-based coding agents, autonomous
software agents, and memory techniques in AI agents: five systematic reviews or mapping
studies and one memory-focused survey.
For each review we extracted: the full strings (when available), the term sets in the agent
block and the memory block, and the representative systems catalogued.
In parallel, a list of 16 systems known from current literature was compiled as recall
test cases, and ten alternative term formulations were tested individually in exploratory
queries on Scopus and arXiv.
\autoref{tab:auditoria-slrs} summarises the reviews consulted and each one's contribution
to the block refinement.

\begin{table}[htbp]
\centering
\caption{Systematic reviews audited to calibrate the search strings.}
\label{tab:auditoria-slrs}
\small
\begin{tabular}{p{3.2cm}p{1cm}p{2.2cm}p{1cm}p{4.5cm}}
\toprule
\textbf{Review} & \textbf{Year} & \textbf{Venue} & \textbf{N} & \textbf{Contribution to refinement} \\
\midrule
\citet{hou2024llmse} & 2024 & TOSEM & 395 & Terms in the \textit{coding}/\textit{software engineering} block; venue coverage. \\
\citet{liu2024llmagents} & 2024 & TOSEM & 117 & Vocabulary for autonomous software agents. \\
\citet{zhang2024unifying} & 2024 & SCIS & 52 & Planning and memory terms in LLM agents. \\
\citet{jin2024llmagents} & 2024 & arXiv & 90 & Variants of \textit{LLM-based agent}; synonyms for block C1. \\
\citet{wang2024llmagents} & 2024 & ASE & 106 & Confirmation of the term \textit{language agent}; role vocabulary. \\
\citet{du2026memory} & 2026 & arXiv & --- & Memory taxonomy for LLM agents; vocabulary for block C2. \\
\bottomrule
\end{tabular}
\end{table}

The audit identified three gaps in the initial string draft: absence of variants without the
\textit{-ing} suffix (\textit{code assistant}), absence of the emerging term \textit{agentic
coding} (already adopted as a venue category in 2025), and absence of the \textit{SWE agent}
form as a proper-name class label.
The corresponding gaps in the memory block covered \textit{persistent memory},
\textit{persistent context}, and \textit{cross-task}, which recur in the seed papers and in
the audited reviews.

\subsection{Seed-paper recovery (Quasi-Gold Standard)}
\label{ap:qgs}

Quasi-Gold Standard validation consists in checking whether a search strategy recovers a
pre-selected set of papers representative of the scope, identified independently of the
strategy being validated~\citep{mourao2020hybrid}.
Seven seed papers were defined, chosen for covering different persistence architectures
(experience bank, episodic memory, commit-oriented memory, \textit{git}-inspired context)
and for including at least one negative result.
The strategy underwent two refinement cycles documented in the versioned protocol, raising
recovery from four to six of the seven seed papers; the remaining case (CodeMEM) was not
captured by any database query and was added manually as a protocol amendment, as
described in the Method section.
\autoref{tab:qgs-recuperacao} consolidates the per-database recovery status.

\begin{table}[htbp]
\centering
\caption{Seed-paper recovery by database and method. \textit{String} indicates recovery via
the primary Boolean search; \textit{Manual} indicates a documented protocol amendment
for a seed paper not captured by any query.}
\label{tab:qgs-recuperacao}
\small
\begin{tabular}{p{3.5cm}p{1cm}p{1.5cm}p{1.2cm}p{1.5cm}p{2cm}}
\toprule
\textbf{Seed paper} & \textbf{Year} & \textbf{Scopus} & \textbf{arXiv} & \textbf{S2 ID} & \textbf{Recovered} \\
\midrule
Du (memory survey)~\citep{du2026memory} & 2026 & no & yes & yes & Yes (string) \\
Joshi et al.\ (SWE-Bench-CL)~\citep{joshi2025swebenchcl} & 2025 & no & yes & yes & Yes (string) \\
Wang et al.\ (CodeMEM)~\citep{wang2026codemem} & 2026 & no & no & yes & Yes (Manual) \\
Deng et al.\ (MemCoder)~\citep{deng2026memcoder} & 2026 & no & yes & yes & Yes (string) \\
Wu et al.\ (GCC)~\citep{wu2025gcc} & 2025 & no & yes & yes & Yes (string) \\
Wang et al.\ (Repository Memory)~\citep{wang2026improving} & 2026 & no & yes & yes & Yes (string) \\
Lindenbauer et al.\ (CTIM-Rover)~\citep{lindenbauer2025ctimrover} & 2025 & no & yes & yes & Yes (string) \\
\bottomrule
\end{tabular}
\end{table}

The complete absence of the seed papers from Scopus reflects a characteristic of the field:
most analysed systems are 2025--2026 \textit{preprints} still outside the indexing cycle of
traditional venues.
This diagnosis was confirmed by auxiliary Scopus queries (searches on system proper names),
which returned zero results, and justifies the centrality of citation chasing as a third
PRISMA method in the adopted strategy.

\subsection{PRISMA-S compliance}
\label{ap:prisma-s}

\autoref{tab:prisma-s} presents the report of the search strategy following the 16 items of
the PRISMA-S extension~\citep{rethlefsen2021prismas} to PRISMA
2020~\citep{page2021prisma}.
The items cover source naming, platforms, limits, supplementary search methods, reference
management, and updates.

\begin{longtable}{p{0.5cm}p{4cm}p{9cm}}
\caption{PRISMA-S report (16 items) of the search strategy.}
\label{tab:prisma-s} \\
\toprule
\textbf{\#} & \textbf{Item} & \textbf{Description} \\
\midrule
\endfirsthead
\toprule
\textbf{\#} & \textbf{Item} & \textbf{Description} \\
\midrule
\endhead
\bottomrule
\endfoot
1 & Database names & Scopus and arXiv. \\
2 & Platform/interface & Scopus via the Elsevier API (\texttt{scopus-mcp}); arXiv via REST API with full pagination (\texttt{fetch\_arxiv.py}). \\
3 & Study registers & No study registers (for example, PROSPERO) were consulted. \\
4 & Online resources & No organisational websites were consulted separately. \\
5 & Citation chasing & Forward and backward via Semantic Scholar from seven seed papers. \\
6 & Author contact & No author was contacted. \\
7 & Other methods & One documented manual addition: CodeMEM~\citep{wang2026codemem}, a seed paper not captured by the database strings; counted under the \textit{Manual} row of \autoref{tab:search}. See the amendment dated 2026-03-28 in the replication package. \\
8 & Full strategies & Terms by conceptual block in \autoref{tab:search-strings}; exact per-database strings in the replication package (DOI: 10.5281/zenodo.19463131), versioned in \textit{git}. \\
9 & Limits and restrictions & Scopus: \texttt{PUBYEAR > 2022}; arXiv: categories \texttt{cs.SE}, \texttt{cs.AI}, \texttt{cs.CL} and post-search filter \texttt{--year-from 2023}; English and Portuguese languages. \\
10 & Strategy development & Mapping of three conceptual blocks, string construction with OR/AND, validation against seven seed papers, and two iterative refinement cycles documented in the versioned protocol. \\
11 & Peer review & Strategy reviewed by the supervisor before final execution. \\
12 & Record management & Zotero with Better BibTeX for management; \texttt{.bib} files versioned in \textit{git}. \\
13 & Total records & 1,007 from the databases (Scopus: 39; arXiv: 968), 186 unique from citation chasing, and 1 from documented manual addition (CodeMEM); total of 1,194 before deduplication. After removing 83 duplicates, 1,111 records entered title/abstract screening. \\
14 & Deduplication & Elicit Pro deduplicates automatically on import; subsequent manual removal totalled 83 duplicates. \\
15 & Automation tools & Scopus via MCP, arXiv via Python script, assisted screening via Elicit Pro, citation chasing via Semantic Scholar MCP. \\
16 & Updates & Search to be re-run before submission if more than three months have elapsed since the initial execution. \\
\end{longtable}
